\documentclass[11pt,a4paper]{article}

\usepackage[margin=2.5cm]{geometry}
\usepackage{amsmath,amssymb,amsthm}
\usepackage{hyperref}
\usepackage{booktabs}

\theoremstyle{definition}
\newtheorem{definition}{Definition}[section]
\newtheorem{remark}[definition]{Remark}

\theoremstyle{plain}
\newtheorem{proposition}[definition]{Proposition}
\newtheorem{observation}[definition]{Observation}

\newcommand{\R}{\mathbb{R}}
\newcommand{\Z}{\mathbb{Z}}
\newcommand{\Q}{\mathbb{Q}}

\title{\textbf{Dual Meixner--Pollaczek Functions\\
and a Matching Matrix over a Bessel Ring}}
\author{Robin Langer and Claude Code}
\date{August 2026}

\begin{document}
\maketitle

\begin{abstract}
The Meixner--Pollaczek weight $W(x) = |\Gamma(\tfrac14+ix)|^2$ does not
admit a nice square root: $\sqrt{W}$ has branch points at $x = \pm i/4$,
and the resulting matching matrix (companion paper) has no algebraic
structure. But $W$ factors \emph{analytically},
$W = \Gamma(\tfrac14+ix) \cdot \Gamma(\tfrac14-ix)$, giving a
\emph{dual pair} of Meixner--Pollaczek functions
$F_m = \Gamma(\tfrac14+ix) \widehat{P}_m$ and
$G_m = \Gamma(\tfrac14-ix) \widehat{P}_m$, biorthonormal under the
bilinear pairing $\int fg\,dx$. Projecting $F_m$ onto the polynomial
basis in $L^2(W)$ yields matching coefficients $c_{mk}$ obeying a
phase law $c_{mk} = (-i)^{|m-k|}\,|c_{mk}|$, and --- the main
observation --- after rescaling by $2^m m!\, 2^k k!$ every entry lies in
the rank-two lattice $\Z A \oplus \Z B$ over the two Bessel constants
\[
  A = \sqrt{2e}\; K_0(\tfrac12), \qquad
  B = \sqrt{2e}\; K_1(\tfrac12).
\]
The integrality lost with $\sqrt{W}$ is recovered by the dual pair,
with two transcendentals replacing the Hermite matrix's single constant
$\sqrt{2\pi/3}$. The integer pairs grow while the entries shrink:
each row furnishes rational approximations to
$K_1(\tfrac12)/K_0(\tfrac12)$.
\end{abstract}

%% ===================================================================
\section{The Dual Pair}
%% ===================================================================

Let $P_n = P_n^{(1/4)}(\,\cdot\,; \tfrac{\pi}{2})$ be the
Meixner--Pollaczek polynomials, orthogonal on $\R$ with respect to
$W(x) = |\Gamma(\tfrac14+ix)|^2$, and $\widehat{P}_n = P_n/\sqrt{h_n}$
the orthonormal family, $h_0 = \int W = \sqrt2\,\pi^{3/2}$.

The Gaussian weight factors as
$e^{-t^2} = e^{-t^2/2} \cdot e^{-t^2/2}$: two equal positive halves,
which is why a single orthonormal basis (the Hermite functions)
suffices there. The Meixner--Pollaczek weight refuses this splitting
--- $\sqrt{W}$ has branch points at $\pm i/4$ --- but factors into an
analytic function and its reflection:
\[
  W(x) = \Gamma(\tfrac14 + ix)\cdot\Gamma(\tfrac14 - ix).
\]

\begin{definition}
The \emph{dual Meixner--Pollaczek functions} are
\[
  F_m(x) = \Gamma(\tfrac14 + ix)\, \widehat{P}_m(x),
  \qquad
  G_m(x) = \Gamma(\tfrac14 - ix)\, \widehat{P}_m(x).
\]
\end{definition}

$F_m$ is meromorphic with poles only at $x = i(\tfrac14 + n)$,
$n \in \mathbb{N}$ (upper half-plane); $G_m$ has the reflected poles.
Under the \emph{bilinear} pairing $\langle f, g\rangle = \int_\R fg\,dx$
(no conjugate),
\[
  \langle F_m, G_k \rangle
  = \int_\R \widehat{P}_m \widehat{P}_k\, W\, dx = \delta_{mk}:
\]
a biorthonormal system, with each family carrying one \emph{whole}
$\Gamma$-factor --- never a half.

\begin{remark}
$F_m = e^{i\theta(x)}\, \Phi_m$ where
$\Phi_m = \widehat{P}_m\sqrt{W}$ is the ``obvious''
Meixner--Pollaczek function and
$e^{i\theta(x)} = \Gamma(\tfrac14+ix)/|\Gamma(\tfrac14+ix)|$ is a
unimodular phase (the $\Gamma$-part of the Riemann--Siegel theta
function). Multiplication by $e^{i\theta}$ is unitary on $L^2(\R)$, so
$\{F_m\}$ is also orthonormal in the usual sesquilinear sense, and
$G_m = \overline{F_m}$. By Bump--Ng, $F_m$ is (up to constants) the
Mellin transform of the even Hermite function $\psi_{2m}$ on the
critical line.
\end{remark}

%% ===================================================================
\section{The Matching Coefficients}
%% ===================================================================

Projecting $F_m$ onto the polynomial basis in $L^2(W)$:
\begin{equation}\label{eq:cmk}
  c_{mk} = \langle F_m, \widehat{P}_k \rangle_W
  = \int_\R \widehat{P}_m \widehat{P}_k\;
    \Gamma(\tfrac14+ix)^2\, \Gamma(\tfrac14-ix)\, dx.
\end{equation}
The coefficients of $G_m$ are the complex conjugates. Three whole
$\Gamma$-factors appear --- no branch points anywhere.

\begin{proposition}[Checkerboard]
$c_{mk} = c_{km}$, and $x \mapsto -x$ gives
$c_{mk} = (-1)^{m+k}\, \overline{c_{mk}}$: entries with $m+k$ even are
real, entries with $m+k$ odd are purely imaginary. There is no parity
vanishing --- the matrix is full.
\end{proposition}

\begin{observation}[Phase law]\label{obs:phase}
For all computed entries ($m, k \le 11$),
\[
  c_{mk} = (-i)^{\,|m-k|}\; |c_{mk}|.
\]
The phase depends only on $|m-k|$ --- an echo of the Fourier
eigenvalues $(-i)^m$ of the Hermite functions on the Mellin side.
\end{observation}

Table~\ref{tab:numeric} gives the upper-left $10 \times 10$ corner
(\texttt{mp\_dual\_matching.py}, 40-digit quadrature, cache
\texttt{mp\_dual\_matching.pkl}).

\begin{table}[ht]
\centering
\caption{$c_{mk}$ for $0 \le m, k \le 9$. Real for $m+k$ even,
imaginary for $m+k$ odd; the phase is $(-i)^{|m-k|}$.}
\label{tab:numeric}
\scriptsize
\setlength{\tabcolsep}{3.5pt}
\begin{tabular}{r|rrrrrrrrrr}
\toprule
$m \backslash k$ & $0$ & $1$ & $2$ & $3$ & $4$ & $5$ & $6$ & $7$ & $8$ & $9$ \\
\midrule
$0$ & $2.155$ & $-.890i$ & $-.659$ & $.504i$ & $.425$ & $-.365i$ & $-.324$ & $.291i$ & $.265$ & $-.244i$ \\
$1$ & $-.890i$ & $.542$ & $-.226i$ & $-.301$ & $.222i$ & $.223$ & $-.190i$ & $-.182$ & $.165i$ & $.156$ \\
$2$ & $-.659$ & $-.226i$ & $.335$ & $-.084i$ & $-.213$ & $.125i$ & $.162$ & $-.124i$ & $-.135$ & $.116i$ \\
$3$ & $.504i$ & $-.301$ & $-.084i$ & $.235$ & $-.039i$ & $-.169$ & $.076i$ & $.133$ & $-.086i$ & $-.111$ \\
$4$ & $.425$ & $.222i$ & $-.213$ & $-.039i$ & $.175$ & $-.024i$ & $-.139$ & $.051i$ & $.114$ & $-.062i$ \\
$5$ & $-.365i$ & $.223$ & $.125i$ & $-.169$ & $-.024i$ & $.140$ & $-.017i$ & $-.117$ & $.037i$ & $.099$ \\
$6$ & $-.324$ & $-.190i$ & $.162$ & $.076i$ & $-.139$ & $-.017i$ & $.118$ & $-.012i$ & $-.101$ & $.028i$ \\
$7$ & $.291i$ & $-.182$ & $-.124i$ & $.133$ & $.051i$ & $-.117$ & $-.012i$ & $.102$ & $-.009i$ & $-.089$ \\
$8$ & $.265$ & $.165i$ & $-.135$ & $-.086i$ & $.114$ & $.037i$ & $-.101$ & $-.009i$ & $.089$ & $-.007i$ \\
$9$ & $-.244i$ & $.156$ & $.116i$ & $-.111$ & $-.062i$ & $.099$ & $.028i$ & $-.089$ & $-.007i$ & $.080$ \\
\bottomrule
\end{tabular}
\end{table}

%% ===================================================================
\section{Closed Forms: the Bessel Ring}
%% ===================================================================

The base integral has a closed form. The three-$\Gamma$ Mellin--Barnes
integral evaluates in terms of Tricomi's confluent hypergeometric
function,
\[
  \int_\R \Gamma(a+ix)\Gamma(b+ix)\Gamma(c-ix)\, dx
  = 2\pi\, \Gamma(a+c)\Gamma(b+c)\; U(a+c,\, a-b+1,\, 1),
\]
and $U(\tfrac12, 1, z) = e^{z/2} K_0(z/2)/\sqrt{\pi}$ gives, at
$a = b = c = \tfrac14$,
\[
  \int_\R \Gamma(\tfrac14+ix)^2\,\Gamma(\tfrac14-ix)\, dx
  = 2\pi^{3/2} \sqrt{e}\; K_0(\tfrac12)
  = 16.97348192607767\ldots
\]
(verified to 20 digits). Dividing by $h_0 = \sqrt2\,\pi^{3/2}$:
\[
  c_{00} = \sqrt{2e}\; K_0(\tfrac12).
\]

Since $(a+ix)\Gamma(a+ix) = \Gamma(a+1+ix)$, every polynomial moment
of the kernel reduces to three-$\Gamma$ integrals with
integer-shifted parameters, hence to \emph{contiguous} $U$-values ---
all expressible over the two generators $U(\tfrac12,1,1)$ and
$U(\tfrac32,1,1)$, i.e.\ over
\[
  A = \sqrt{2e}\, K_0(\tfrac12) = 2.155416\ldots, \qquad
  B = \sqrt{2e}\, K_1(\tfrac12) = 3.862231\ldots
\]

\begin{observation}[Integrality]\label{obs:integrality}
Define the rescaled raw entries
\[
  T_{mk} = \frac{2^m m!\; 2^k k!}{h_0}
  \int_\R P_m P_k\, \Gamma(\tfrac14+ix)^2\,\Gamma(\tfrac14-ix)\, dx.
\]
Then (dropping the factor $i$ when $m+k$ is odd) every computed entry
satisfies
\[
  T_{mk} \in \Z A \oplus \Z B,
\]
confirmed by PSLQ at 40 digits for all $m \le 3$, $k \le 5$ except
$(3,5)$, where the integer coefficients exceed the search bound.
\end{observation}

\begin{table}[ht]
\centering
\caption{Exact entries $T_{mk} = \mathfrak{a}_{mk} A +
\mathfrak{b}_{mk} B$ recovered by PSLQ.}
\label{tab:exact}
\small
\begin{tabular}{cl@{\qquad}cl}
\toprule
$(m,k)$ & $T_{mk}$ & $(m,k)$ & $T_{mk}$ \\
\midrule
$(0,0)$ & $A$                   & $(1,3)$ & $2575A - 1440B$ \\
$(0,1)$ & $3A - 2B$             & $(1,4)$ & $39411A - 21978B$ \\
$(0,2)$ & $12B - 23A$           & $(1,5)$ & $394932B - 707389A$ \\
$(0,3)$ & $138B - 241A$         & $(2,2)$ & $2369A - 1320B$ \\
$(0,4)$ & $3265A - 1800B$       & $(2,3)$ & $33743A - 18834B$ \\
$(0,5)$ & $52907A - 29706B$     & $(2,4)$ & $317892B - 569719A$ \\
$(1,1)$ & $12B - 21A$           & $(2,5)$ & $6190602B - 11092229A$ \\
$(1,2)$ & $114B - 205A$         & $(3,3)$ & $295572B - 529549A$ \\
        &                       & $(3,4)$ & $5399634B - 9675553A$ \\
\bottomrule
\end{tabular}
\end{table}

\begin{remark}[Diophantine near-cancellation]
The integer pairs grow rapidly while $T_{mk} \to 0$, so each entry is a
near-cancellation: $\mathfrak{b}_{mk}/\mathfrak{a}_{mk} \approx -A/B$,
i.e.\ the rows generate rational approximations to
\[
  \frac{B}{A} = \frac{K_1(\tfrac12)}{K_0(\tfrac12)} = 1.791825\ldots
  \qquad \Bigl(\text{e.g. } \tfrac{52907}{29706} = 1.78103\ldots\Bigr).
\]
This is the Pad\'e/continued-fraction mechanism behind irrationality
results for Bessel ratios (the Gauss continued fraction), here arising
spontaneously from an approximation problem in $L^2(W)$.
\end{remark}

%% ===================================================================
\section{Comparison with the Self-Dual Matrix}
%% ===================================================================

Against the $\sqrt{W}$ matching matrix of the companion paper:

\begin{center}
\small
\setlength{\tabcolsep}{4pt}
\begin{tabular}{lll}
\toprule
 & self-dual ($\sqrt{W}$) & dual pair ($\Gamma^{+}, \Gamma^{-}$) \\
\midrule
$\Gamma$-count in integrals & $3/2$ (branch points) & $3$ (poles only) \\
parity & $b_{mk} = 0$, $m+k$ odd & full matrix, real/imaginary \\
algebraic structure & none (PSLQ negative) & $\Z A \oplus \Z B$ \\
base constant & $20.9154\ldots$ (no closed form) & $2\pi^{3/2}\sqrt{e}\,K_0(\tfrac12)$ \\
biorthogonal expansion & --- & $\langle F_m, G_k\rangle = \delta_{mk}$ (exact) \\
\bottomrule
\end{tabular}
\end{center}

The decay of $|c_{0k}|$ remains slow (the pole of $F_m$ at $x = i/4$
is still at finite distance; local log-log slope $\approx 0.7$ at
$k = 11$), so the dual pair does not speed up $L^2$ truncation. What it
changes is the \emph{arithmetic}: the coefficients are now exact
elements of a rank-two Bessel lattice, computable by residue calculus
rather than quadrature, and the biorthogonal expansion itself
($F$ against $G$) is exact at every order.

\sloppy
In the Hermite case a single transcendental $\sqrt{2\pi/3}$ carried an
integer matrix with a Wick matching interpretation, and the modified
Bessel function $I_0$ appeared in the diagonal's generating function.
Here \emph{two} transcendentals $K_0(\tfrac12)$ and $K_1(\tfrac12)$
carry a pair of integer matrices
$(\mathfrak{a}_{mk}, \mathfrak{b}_{mk})$.
Open: a combinatorial law for $(\mathfrak{a}, \mathfrak{b})$ (the
denominators $2^m m!\, 2^k k!$ and the diagonal pattern suggest a
weighted-path model on the Jacobi matrix), a proof of the phase law
(Observation~\ref{obs:phase}), and the exact decay exponent.

\end{document}
